\documentclass[12pt,a4paper]{article}

% ============================================================
% NYCU / Arete Honors Program Project Paper Template
% Compiler: XeLaTeX
% 使用方式：
% 1. 修改「基本資料」區塊。
% 2. 替換摘要、關鍵字、各章正文與參考文獻。
% 3. Overleaf 請使用 XeLaTeX，並至少 Recompile 兩次以更新目錄。
% ============================================================

\usepackage[a4paper,top=2.5cm,bottom=2.5cm,left=2.5cm,right=2.5cm,footskip=1cm]{geometry}
\usepackage{fontspec}
\usepackage{xeCJK}
\usepackage{array}
\usepackage{setspace}
\usepackage{indentfirst}
\usepackage{cite}
\usepackage[hidelinks]{hyperref}
\usepackage{zhnumber}
\usepackage{titlesec}
\usepackage{tocloft}
\usepackage{booktabs}
\usepackage{graphicx}
\usepackage{tikz}
\usepackage{enumitem}
\usepackage{caption}

\usetikzlibrary{arrows.meta,positioning}

% ---------- 字體設定 ----------
% Overleaf 若缺少 Noto Serif CJK TC，可改成 AR PL UMing TW 或其他可用中文字體。
\setCJKmainfont{Noto Serif CJK TC}
\setmainfont{TeX Gyre Termes}

% ---------- 基本版面 ----------
\linespread{1.3}
\pagestyle{empty}
\setlength{\parindent}{2em}
\setlength{\parskip}{0pt}
\setlist[itemize]{leftmargin=2em}
\setlist[enumerate]{leftmargin=2em}
\captionsetup{font=small,labelfont=bf}

% ---------- 章節與目錄 ----------
\renewcommand{\thesection}{第\zhnum{section}章\hspace{0.5em}}
\renewcommand{\thesubsection}{\arabic{section}.\arabic{subsection}}
\titleformat{\section}{\Large\bfseries}{\thesection}{0em}{}
\titleformat{\subsection}{\large\bfseries}{\thesubsection}{0.75em}{}
\setlength{\cftsecnumwidth}{4.5em}
\setlength{\cftsubsecnumwidth}{3em}

\renewcommand{\contentsname}{\hfill 目錄\hfill\null}
\renewcommand{\listfigurename}{\hfill 圖目錄\hfill\null}
\renewcommand{\listtablename}{\hfill 表目錄\hfill\null}
\renewcommand{\figurename}{圖}
\renewcommand{\tablename}{表}
\renewcommand{\refname}{參考文獻}

% ============================================================
% 基本資料：之後使用模板時主要改這裡
% ============================================================
\newcommand{\schoolzh}{國立陽明交通大學}
\newcommand{\programzh}{博雅書苑百川學士學位學程}
\newcommand{\coursezh}{專題探索（一）}
\newcommand{\programen}{Arete Honors Program}
\newcommand{\schoolen}{National Yang Ming Chiao Tung University}
\newcommand{\courseen}{Project Exploration (I)}

\newcommand{\papertitlelineone}{Your English Paper Title Line One}
\newcommand{\papertitlelinetwo}{Your English Paper Title Line Two}

\newcommand{\studentname}{學生姓名}
\newcommand{\advisorname}{指導教授姓名}
\newcommand{\semester}{2026 Spring}

\begin{document}

% ============================================================
% 封面
% ============================================================
\begin{titlepage}
  \thispagestyle{empty}
  \begin{center}
    \vspace{0.6cm}

    \fontsize{18}{22}\selectfont \schoolzh \\
    \fontsize{18}{22}\selectfont \programzh \\
    \fontsize{18}{22}\selectfont \coursezh \\

    \vspace{1cm}

    \fontsize{14}{17}\selectfont \programen \\
    \fontsize{16}{19}\selectfont \schoolen \\
    \fontsize{16}{19}\selectfont \courseen \\

    \vspace{3.1cm}

    \fontsize{18}{22}\selectfont \papertitlelineone \\
    \fontsize{18}{22}\selectfont \papertitlelinetwo \\

    \vspace{3.1cm}

    \begin{tabular}{@{}r@{\hspace{0.8em}}c@{\hspace{0.8em}}l@{}}
    \makebox[4em][s]{學\hfill 生} & ： & \studentname \\
    \makebox[4em][s]{指導教授} & ： & \advisorname \\
    \end{tabular}

    \vfill

    \fontsize{18}{22}\selectfont \semester\\
  \end{center}
\end{titlepage}

% ============================================================
% 前置頁：中文摘要、目錄、圖目錄、表目錄
% ============================================================
\pagestyle{plain}
\pagenumbering{roman}
\setcounter{page}{1}

\clearpage
\phantomsection
\addcontentsline{toc}{section}{中文摘要}
\section*{\centering 中文摘要}
請在此撰寫中文摘要。摘要宜簡要說明研究背景、研究問題、研究方法、主要結果與貢獻。若依正式論文規範，可控制在一頁以內，並在摘要後列出 5 至 7 個關鍵字。

\noindent\textbf{關鍵字：}關鍵字一、關鍵字二、關鍵字三、關鍵字四、關鍵字五

\newpage
\tableofcontents

\newpage
\addcontentsline{toc}{section}{圖目錄}
\listoffigures

\newpage
\addcontentsline{toc}{section}{表目錄}
\listoftables

% ============================================================
% 正文：從第一章開始使用阿拉伯數字頁碼
% ============================================================
\newpage
\pagenumbering{arabic}
\setcounter{page}{1}

\section{緒論}

\subsection{研究背景}
請說明研究主題的背景、重要性，以及目前問題所在。此段可包含領域現況、實務需求、學術缺口或課程脈絡。

\subsection{研究問題}
請明確列出本研究欲回答的核心問題。可使用一段文字描述，也可使用條列式呈現。

\subsection{研究目的與貢獻}
請說明本研究的目的與預期貢獻。例如：
\begin{itemize}
  \item 貢獻一。
  \item 貢獻二。
  \item 貢獻三。
\end{itemize}

\section{文獻探討}

\subsection{相關研究一}
請整理與本研究直接相關的文獻，並指出其方法、發現與限制。例如可使用引用格式 \cite{example2026paper}。

\subsection{相關研究二}
請比較不同研究之間的差異，並說明本研究如何延伸或補足既有工作。

\section{研究方法}

\subsection{研究架構}
請說明整體研究流程。若需要流程圖，可參考圖 \ref{fig:pipeline-example}。

\begin{figure}[htbp]
  \centering
  \begin{tikzpicture}[
    node distance=1.5cm,
    box/.style={draw, rounded corners, align=center, minimum width=3.2cm, minimum height=1cm},
    arrow/.style={-Latex, thick}
  ]
    \node[box] (input) {輸入資料};
    \node[box, right=of input] (method) {研究方法};
    \node[box, right=of method] (output) {分析結果};
    \draw[arrow] (input) -- (method);
    \draw[arrow] (method) -- (output);
  \end{tikzpicture}
  \caption{研究流程示意圖}
  \label{fig:pipeline-example}
\end{figure}

\subsection{資料與工具}
請描述使用的資料、工具、模型、實驗環境或分析方法。

\subsection{實作或分析流程}
請詳細說明實作步驟、演算法流程、系統架構或分析程序。

\section{實驗結果與分析}

\subsection{結果呈現}
請呈現主要結果。若需要表格，可參考表 \ref{tab:result-example}。

\begin{table}[htbp]
  \centering
  \caption{實驗結果範例表}
  \label{tab:result-example}
  \begin{tabular}{lcc}
    \toprule
    方法 & 指標一 & 指標二 \\
    \midrule
    方法 A & 0.80 & 0.75 \\
    方法 B & 0.85 & 0.78 \\
    \bottomrule
  \end{tabular}
\end{table}

\subsection{結果討論}
請解釋結果代表的意義，並說明是否符合研究預期。

\section{討論}

\subsection{研究發現}
請整理本研究的主要發現，並連結回研究問題。

\subsection{限制}
請說明資料、方法、實作、評估或理論上的限制。

\section{結論與未來工作}
請總結本研究完成的工作、主要貢獻，以及未來可以延伸的方向。

\clearpage
\addcontentsline{toc}{section}{參考文獻}
\begin{thebibliography}{99}

\bibitem{example2026paper}
Author Name.
\newblock Paper Title.
\newblock \textit{Conference or Journal Name}, Year.

\end{thebibliography}

\end{document}
